\chapter{March}

\section[March 2nd]{March 2\textsuperscript{nd}}
The class started with the story ``Malala's Magic Pencil''. It narrated the life of Malala, a girl who liked to draw and study. After noticing the problems in Pakistan, she started to dream about helping to solve the issues facing her people, such as war and extremism. The story concluded by showing how her writings played an important role in bringing worldwide attention to her people's situation.

I believe the professor chose the story to represent the importance of proper writing and how powerful good texts can be. This relates directly to the current subject of academic writing. Quality academic productions can provide small advancements in human knowledge and science as a whole. Therefore, similar to Malala's writing, good academic writing can drive change and contribute to a better world, even if only by a small amount.

Following the story, there was a round of quick introductions among classmates. Then, we discussed the course plan, reviewing the evaluation process and the topics that will be covered throughout the semester.

It seems the classes will be interesting. There will be a strong emphasis on how different English is from Portuguese in terms of writing. The focus will go beyond simple translations and delve into the logic itself, especially when dealing with summaries.

The idea of improving my writing and seeing this familiar language through new lenses is exciting. Judging by the reading of the children's book, I believe we will read and write quite a bit. We will become familiar with different texts and learn good argumentation and sentence construction from a highly technical and objective perspective.

The slides for this class can be accessed in the appendix, section \ref{app:mar02}.

\section[March 9th]{March 9\textsuperscript{th}}

The class started with the book ``I Am Me: A Book of Authenticity''. It is a book about the importance of being yourself and understanding how everyone is unique in their own way, with their talents, styles, and imperfections. It is mostly for self-confidence. I believe that such readings are more for my classmates who teach children, helping them with reading texts for younger audiences. The final tip discussing mindful meditation was more general. The whole story prior was unnecessary to teach that concept. This makes me think that those story introductions are more of an extension of the ``English for Kids'' classes.

Regarding the meditation itself, it was not particularly useful for me. Repetitions about concentrating only on the good things about ourselves fail to show the full picture. The actual reality can be more nuanced. We are capable of having genuine flaws and committing truly evil actions. Understanding that reality and making an effort to improve ourselves is necessary to become a better person.

The next topic of discussion was about the concept of ``research''. We need to search for topics systematically, analyze the data, and so on. In scientific research, we can often analyze and review other published works and collections without needing to produce the data ourselves. Then we talked about the concept of ``scientific''. It refers to methods, principles, or knowledge based on systematic observation, experimentation, and evidence. Modern science uses the scientific method. It is important to follow a logical approach, to be reproducible, and to maintain a strict structure to draw valid conclusions.

We learned more about scientific life, such as the Lattes curricula and ORCID. It is important to be able to publish works, get recognition, and strengthen one's name in academia. Regarding the course, we will follow the cycle of reading, writing, and gathering feedback so we can read and write again.

\section[March 16th]{March 16\textsuperscript{th}}

The assignment for this day was to read the first chapter of the book ``Produção textual acadêmica'' by Lêda Maria Braga Tomitch. The book covers the general concepts of reading, learning, and writing. The introduction to reading was rather interesting because it goes beyond the mere concept of skimming over words and understanding their meaning in the context of phrases. The text brings about how there are many different types of text and of reading, each with its own form of understanding. We can skim over a magazine for entertainment while waiting in the dentist's office, or sit down and read the textbook to prepare for an exam. Reading involves the reader, the writer, and the text itself.

Learning is the process of gathering new information and storing it in such a manner that it can be accessed. Such a definition can be seen as more or less dynamic, since instead of simply ``piling up'' knowledge, our own knowledge base actually interacts and changes with the new information that is being acquired.

Writing is the process of gathering our thoughts and passing them to another person through text; it is a social process because the knowledge and understanding of the reader must be taken into consideration. It is also a great form of learning and using the information acquired, because when writing a text, we are forced to organize our thoughts and make them as clear as possible to the reader, proving a certain mastery of the stored knowledge.

The textbook provides a good introduction to the class, showing more of the ``inner workings'' of reading, learning, and writing, making me curious to see how we will apply such concepts during class.

The second assignment was to write a summary about an academic text called ``Fluency Revisited''. Writing such a summary was quite interesting in order for me to practice my writing. I could notice that keeping good grammar throughout a long text summarizing ideas can be way more challenging than I initially thought.

The actual class started with another children's book called ``I am Human: A Book of Empathy'' by Susan Verde. It is once again another book about the idea that humans are not perfect, but are not special. The book complements the previous one, talking about how we can hurt others, have fear, be timid, have choices, and so on. At the end, the book suggests that we should orient our actions with thoughtfulness, kindness, fairness, forgiveness, and so on. Also, the book invited the readers to hope and dream. The professor reflects about how we should learn to be the best version of ourselves, especially comparing ourselves with ourselves and not others.

The main thought that goes through my mind during those reading sessions is ``How is this necessary?''. The reading session from today lasted until 7:45 pm. The class started at 6:30 pm, and the reading session takes more than an hour, going through each page describing the illustrations in detail; everything gets really repetitive and tiring. And the messages and reflections that the book brings are shallow and generic.

We are invited to act with fairness, but what is fairness? We should value ourselves, but to what extent? When does self-respect end and pride begin? General common sense messages such as those work fine with children, but in an academic setting, I expected something more profound or better argumentation and thoughts. If I were to start with important reflections about what it means to be good, have empathy, and so on, the class would benefit from deeper texts, such as parts from Aristotle's Nicomachean Ethics, Saint Thomas Aquinas' Summa Theologiae, Kant's The Metaphysics of Morals, or other well-argued texts about love, compassion, happiness, peace, social life, human nature, forgiveness, the meaning of life, and so on. With such texts, the training of good argumentation and reflection would help to improve our own writing and thinking capabilities beyond the self-help/common sense from children's books.

It is quite visible that we take much time to cover such little ground. That becomes quite frustrating; even though most of the students in this class are teachers, such long readings might be a little exaggerated to teach them to read for children, especially considering that this class is Academic Writing and not English for Kids anymore.

At 8:00 pm, we finally started actually analyzing the ``Fluency Revisited'' text, understanding the distinct aspects of the text that define it as academic, such as keywords, citations, the reference list, details about the publication, use of arguments, and general structure with an introduction featuring important definitions and a brief explanation of what the text will be about, the general development of the text, and a summarized conclusion.

The importance of science is discussed, highlighting how the idea of a ``scientist'' goes beyond the stereotype of white coats and hard sciences. Instead, it involves following the scientific process, manuscript (when the paper is not yet published) or article writing, and publishing, including scientific peer review through anonymous analysis of the manuscripts before publication.

Then, we went through Tomitch's text together in class. We reflected upon the different types of knowledge that can be acquired through reading. These are: content (the actual ideas being expressed in the text), formal (the general structure on how those ideas are laid out, like in the format of introduction, definitions, argumentation, conclusion, and so on), and linguistic (the general vocabulary, grammar, and relations between the words themselves).

The author asks why students need to read so much and write so much. The answer brought up in class is that we need reading to get and store ideas, and also to learn about formal and linguistic knowledge. Writing helps with organizing and actively using the stored information.

Professor Tomitch used a big diagram about reading that can be seen in the appendix, section \ref{app:tomitch}. The diagram shows the different processes that happen during reading, like decoding, literal comprehension, inferential comprehension, and comprehension monitoring.

I learned about the concept of schemas, like a ``code of conduct'', a knowledge that we have about how to behave and the general idea of how to behave and interact with different environments in different contexts.

Also, a new concept is that of procedural knowledge. In the example of riding a bike or driving, we follow procedures; we know how to follow a sequence of actions, such as using the different pedals. Learning it is a difficult process, but as time goes on, those difficult and overwhelming processes start to get more and more automatic until they become second nature. This sequence of steps that we need to take in order to reach the objectives constitutes procedural knowledge.

In decoding, we have two processes: matching and recoding. Decoding is like ``decoding the message from the text'', putting the letters together to form words, and so on.

Literal comprehension is the comprehension of the literal words \textit{ipsis litteris}, the actual literal meaning of the word in itself. Lexical access is when we access the meaning of the word alone. Parsing happens when we understand the word within the syntactical context, like if it is a subject, verb, object, and so on.

Inferential comprehension is a higher-level process; it refers to the moment where we make inferences, and we get meanings beyond what is explicitly laid out in the text. In there, we have integration, summarization, and elaboration (e.g., creating a mind map, integrating the text with previous knowledge to elaborate something new).

Comprehension monitoring does not happen only at the end. Points like ``setting a goal'' and ``selecting strategies'' happen at the beginning of the reading process. Then, we can check if the goal of the reading is reached; if not, we can remediate the strategy and improve it.

Breaking down the reading process into those parts provides a new way of understanding what we do every day; it's like taking a peek into how our mind works and how we interact with texts. This can help in improving our self-assessment of our reading skills and facilitate the process of finding problems with our reading and improving it.

We reviewed some concepts, like ``e.g.'', which stands for exempli gratia from Latin. Acronyms form a word, like FURB, NASA, and Celesc. Abbreviations don't form words, like ABNT.

\section[March 23rd]{March 23\textsuperscript{rd}}
Today I got sick and couldn't attend class. Nevertheless, I read the slides and the text that was covered in class. I also read the story of the day, ``Maybe'' by Kobi Yamada. This is a beautiful book; it reads more like poetry. It develops heavily on the concepts of hope and ``potential''. According to Aristotle, potential is the capacity to become something, and it is a very important concept in philosophy. By definition, potency contrasts with act, which is the actualization of the potential. This makes me reflect on the deeper meaning of human potency: the possibilities feel unlimited, yet at the same time, they are limited by our own nature. I believe it is always important to dream and have high expectations in life, while simultaneously being smart and analytic about our surroundings in order to make the best use of the opportunities that life brings. I never thought that I would be going to Germany again so soon to study a topic that I was never particularly interested in (Project Engineering), but that is where life appears to be guiding me. By following these opportunities, I can make adjustments to my own plans and dreams, actualizing the potential that I truly have.

Moving to some general discussions about Tomitch's book, this class focused on study methods. It is amazing how many methods and nuances there are when deciding how to proceed with studying. The act of actually studying and absorbing content is very challenging. I became particularly good at studying for tests; my grades were always high. My method is quite simple. I just read the text, and during my reading, I stop and rehearse it to myself to be sure that I understood it, then continue until the end. After a while, I return to the text, and without reading its content, I try to explain to myself what it says; if I am happy with the results, I continue. Usually, I like to make my study sessions as stress-free as possible, so I tend to avoid note-taking or more complex constructions. Organizing the ideas inside my brain and giving lectures to myself on the topic tend to work well for me. Naturally, once the test is over, my memory of the topic starts to fade, but never completely. If I need to return to that topic many years later, I feel like I will have a much easier time than the first time. Another topic that I like a lot is exercises. Doing exercise lists helps me step out of my comfort zone. Sometimes I have an idea of what is most important or what a topic presents, but exercises created by third parties or by the authors can throw a curveball at that notion, helping me see the same content in new ways.

Regarding the suggestions in the text and in class, I absolutely despise mind maps. To me, they always felt highly inefficient; they look more like fancy bulleted lists nested into each other. The part that I do not like is the ``fancy'' aspect. It always feels like I spend more time organizing the map and strugglinh making everything fit than actually working with the topic itself. I wasn't very happy with the mind map homework assignment because of that.

Another topic from the text that caught my attention is the SQR3 program: Survey, Question, Read, Recite, and Review. It sounds similar to the way I study. I quickly take a look at what the text has to offer, formulate questions for myself, read to get deeper and answer the problems in my head, stop to recall during the process, and then review everything after it is done to see if everything is in order. I like this method because it can be done lazily. I can grab my cell phone and read the book that I want to study anywhere, and even when I put it away, I can keep going over the steps of Question and Recite inside my head. When I fail to remember or answer something, I simply return to the text.


\section[March 30th]{March 30\textsuperscript{th}}
Today the Professor's daughter got sick, so we had class through Teams even though Professor Bailer had the license to not work; it is admirable that she decided to teach us anyway. Today's story is a recounting of the Alice in Wonderland story. Unfortunately, I have never read or watched the original or any other adaptation, therefore I cannot compare it to the original.

The story appears very childish as does the narration from the Professor with many colors and overly expressive narration. In the book, Alice falls into a hole, shrinks, gets big, talks with hatted caterpillars, the Cheshire Cat appears, goes to a Tea Party, and plays croquet. In the end, it was all a dream. Honestly, this story made me even more confused; the point of the book is a bit confusing, just recount the story in a simple manner with different illustrations? From what my colleagues said, the story remained roughly the same as the original. The good part is that it was quick. Today the story session only lasted 30 minutes, so I don't really mind that it didn't feel very productive. It served as a decent icebreaker to get started with the class.

I should start to read articles from journals that are aligned with my interests and start getting familiarized with them. I am happy that I can choose any topic that I would like to read about and start working with it. Recently, besides my area (Computer Science), I also intend to take a look at productions in Philosophy and Theology. I have interest in these areas and I am curious to check what the current academic production is.

Many sources for academic research were shared; such lists are useful for future research.

Then we went through some important concepts:

\begin{description}
    \item[Phrase] No translation to Portuguese (sometimes compared to Chomsky's \textit{sintagma verbal ou nominal}). A single unit of meaning. E.g., ``The flower''.
    \item[Sentence] \textit{Frase} in Portuguese. A group of words with complete meaning. E.g., ``The flower is yellow.''
    \item[Clause] \textit{Oração} in Portuguese. E.g., ``I like the flower that is yellow.'' (Here, ``that is yellow'' is a subordinate clause).
\end{description}

Then we went through the three main attributes of paragraphs: unity, coherence, and development. This idea of unity is particularly challenging to me; when writing, many ideas come to mind, and making them sound cohesive and choosing just one purpose for the paragraph is really hard.

Important concepts:
\begin{itemize}
    \item \textbf{Topic Sentence:} Sentence that captures the meaning of the paragraph.
    \item \textbf{Supporting Sentences:} Sentences that develop the topic sentence, providing examples, explanations, and so on.
    \item \textbf{Concluding Sentence:} Sentence that revisits the topic sentence and provides a conclusion to the paragraph.
\end{itemize}

In the 14\textsuperscript{th} slide of the class (accessible in Appendix \ref{app:mar30}) we have a comic of Gandhi saying that he understands the secret of life, but he can't figure \textbf{this} out. \textit{This} referring to the paper that he's reading. To me, the joke is that the text he is holding is so poorly written that even someone who understands the secrets of life can't follow it. But the professor translated it giving the sense that Gandhi understands the meaning of life but can't follow it; such interpretation made me quite confused.

When studying a paragraph, we noticed that in English it is acceptable to repeat some words sometimes, especially when the term is really important, because many times synonyms will not convey the same meaning as the one intended. Therefore, a little bit of repetition can be acceptable in order to maintain precision and clarity in the text.

We studied some more examples of paragraphs following the ``hamburger structure''; some examples follow the structure more or less closely.

During class, we were reminded of punctuation marks such as the colon and semi-colon. I found it a bit funny because I use those marks and names constantly in my job, as they are used all the time in programming languages, so I am already quite familiar with all the names of punctuation marks in English due to that.